\documentclass[11pt]{article}
\usepackage{fontspec}
\usepackage[french]{babel}
\usepackage{amsmath}
\usepackage{amssymb}
\usepackage{amsthm}
\usepackage{mathrsfs}
\usepackage{tikz-cd}
\usepackage{enumitem}
\usetikzlibrary{babel}

\newcommand{\printedpage}[1]{\marginpar{\footnotesize p.~#1}}
\newcommand{\msp}[1]{\marginpar{\footnotesize [#1]}}
% dotted arrow as printed in "c : x_0 ....> x_1"
\newcommand{\dotto}{\mathrel{\tikz[baseline=-0.5ex]\draw[dotted,-latex] (0,0)--(1.6em,0);}}

\begin{document}

\printedpage{10}
% head: A. Grothendieck : Dérivateurs, Ch. XII, Édition des Universités de Montpellier II et Paris VII
% foot: Version 20 juin 2003
\textsc{Démonstration}. Soit
\[ F \hookrightarrow U \]
l'inclusion dans $U = \partial X$ de la catégorie-coin \texttt{[}$V =$\texttt{]} $\vcenter{\hbox{\begin{tikzcd}[row sep=small, column sep=small]
0 \arrow[r] \arrow[d] & a \\
b &
\end{tikzcd}}}$ (cf. figure), isomorphe à $\Psi$. On voit aussitôt que c'est une immersion fermée de Dwyer (le plus petit ouvert de $U$ contenant $\Psi$ est justement la catégorie $\tilde{\Psi} = U/t_1$ de tantôt), donc une $W$-cofibration. Si $U \longrightarrow X$ était une $W$-cofibration, il en serait de même du composé $V \longrightarrow X$. On va voir que ce n'est pas le cas.

Soit $\overline{V} = \vcenter{\hbox{\begin{tikzcd}[row sep=small, column sep=small]
0 \arrow[r] \arrow[d] & a \arrow[d] \\
b \arrow[r] & t_1
\end{tikzcd}}}$ l'adhérence de $V$ dans $X$,
\msp{page 14}et considérons l'inclusion
\[ \overline{V} \longrightarrow X, \]
je dis que \texttt{[c]}'est une immersion fermée de Dwyer, i.e. que pour tout $s \in W\backslash\overline{V}$, où $W$ est le plus petit ouvert de $X$ contenant $\overline{V}$, i.e.
\[
W = \quad
\begin{tikzcd}[row sep=small, column sep=small]
0 \arrow[r] \arrow[d] & a \arrow[d] & a' \arrow[l] \arrow[d] \\
b \arrow[r] & {} & a'' \arrow[l] \\
b' \arrow[u] \arrow[r] & b'' \arrow[u] & c, \arrow[l] \arrow[u]
\end{tikzcd}
\]
la catégorie $s\backslash\overline{V}$ a un objet initial. (Les catégories $s\backslash\overline{V}$ qu'on obtient, à partir des cinq choix possibles pour $s$, sont $\boldsymbol{\Delta}^1$ (si $s = a'$ ou $b'$) ou $e$ (pour $s = a''$, $b''$ ou $c$)). Donc $\overline{V} \longrightarrow X$ est une $W$-cofibration.

Ceci posé, on sait qu'il existe, pour [?] l'inclusion (immersion ouverte)
\[ V \hookrightarrow \overline{V} \qquad \text{i.e. } \Psi \hookrightarrow \boldsymbol{\Delta}^1 \times \boldsymbol{\Delta}^1, \]
un foncteur $V \longrightarrow V'$, donnant lieu à un carré cocartésien
\[
\begin{tikzcd}
V \arrow[r, hook, "i"] \arrow[d, "p"'] & \overline{V} \arrow[d, "\overline{p}"] \\
V' \arrow[r, hook, "i'"] & \overline{V'}
\end{tikzcd}
\tag*{$(*)$}
\]
avec
\[ i, p \in W \ \text{(car $V$, $V'$, $\overline{V}$ $W$-asphérique), mais } i', \overline{p} \notin W \ \text{(i.e. $\overline{V'}$ \textit{pas} $W$-asphérique).} \]

Je rappelle qu'on prend d'abord
\printedpage{11}
% head: A. Grothendieck : Dérivateurs, Ch. XII, Édition des Universités de Montpellier II et Paris VII
% foot: Version 20 juin 2003
\[
V'_0 = \left(
\begin{tikzcd}[row sep=small]
 & & a \\
0' \arrow[r, shift left, "\lambda"] \arrow[r, shift right, "\mu"'] & 0 \arrow[ur, "\alpha"] \arrow[dr, "\beta"'] & \\
 & & b
\end{tikzcd}
\right)
\qquad \text{avec} \quad
\left\{
\begin{aligned}
\alpha\lambda &= \alpha\mu \\
\beta\lambda &= \beta\mu
\end{aligned}
\right.
\]
et l'immersion fermée ‘identique’ de $V$ dans $V'$ \texttt{[plutôt} $V'_0$\texttt{]}
\[ p_0 : V \longrightarrow V'_0 \]
\msp{page 15}(qui est telle que
\[ p_{0!}(e_V) = e_{V'_0} \ \ ), \]
puis on prend pour $V'$ la catégorie obtenue en ajoutant un objet final à $V'_0$
\[
V' = \left(
\begin{tikzcd}[row sep=small, column sep=small]
 & & a \arrow[dr, "\rho"] & \\
0' \arrow[r, shift left, "\lambda"] \arrow[r, shift right, "\mu"'] & 0 \arrow[ur, "\alpha"] \arrow[dr, "\beta"'] & & e \\
 & & b \arrow[ur, "\sigma"'] &
\end{tikzcd}
\right)
\qquad \text{avec} \quad
\left\{
\begin{aligned}
\alpha\lambda &= \alpha\mu \\
\beta\lambda &= \beta\mu \\
\rho\alpha &= \sigma\beta
\end{aligned}
\right.
\]
et pour $p$ le foncteur composé
\[
\begin{tikzcd}
V \arrow[r, "p_0"] \arrow[rr, bend right, "p"'] & V'_0 \arrow[r, "p_1"] & V',
\end{tikzcd}
\]
où $p_1$ est l'inclusion (ouverte). On trouve alors le carré $(*)$, avec $\overline{V'}$ non $W$-asphérique, ou ce qui revient au même, $i' \notin W$. Car comme $\overline{V'}\backslash V' \simeq \overline{V}\backslash V$ est la catégorie ponctuelle, $i' \in W$ équivaut à $V'/t$ $W$-asphérique (\footnote[8]{moyennant (W 7 bis).}) (cf. cor. 4, p. 149 dans XII). Or par construction $V'/t \simeq V'_0$, qui n'est pas $W$-asphérique (dans le cas, disons, où $W \subseteq W_1$, car $\pi_1(V'_0) \simeq \mathbf{Z}$).

\msp{page 16}Considérons alors le diagramme de carrés cocartésiens
\[
\begin{tikzcd}
V \arrow[r, hook, "i"] \arrow[d, "p"'] & \overline{V} \arrow[r, hook] \arrow[d, "\overline{p}"] & X \arrow[d, "q"] \\
V' \arrow[r, hook, "i'"] & \overline{V'} \arrow[r, hook] & X'
\end{tikzcd}
\qquad .
\]

\printedpage{12}
% head: A. Grothendieck : Dérivateurs, Ch. XII, Édition des Universités de Montpellier II et Paris VII
% foot: Version 20 juin 2003
Comme $\overline{V} \hookrightarrow X$ est \texttt{[une]} immersion fermée de Dwyer, le carré de droite est $W$-cocartésien, donc $(\overline{p} \in W) \iff (q \in W)$, donc (comme $\overline{p} \notin W$) on a
\[ p \in W, \qquad \text{mais} \qquad q \notin W, \]
ce qui montre bien que $V \longrightarrow X$ n'est pas une $W$-cofibration. Factorisant $V \longrightarrow X$ par $U$ (au lieu de $\overline{V}$), on trouve
\[
\begin{tikzcd}
V \arrow[r] \arrow[d, "p\in W"'] & U \arrow[r] \arrow[d, "q_U\in W"] & X \arrow[d, "q\notin W"] \\
V' \arrow[r] & U' \arrow[r] & X'
\end{tikzcd}
\qquad ,
\]
où $q_U \in W$ puisque $p \in W$ et $V \longrightarrow U$ est \texttt{[une]} immersion de Dwyer. Cela rend explicite que $U \longrightarrow X$ n'est pas une $W$-cofibration, puisque
\[ q_U \in W, \qquad \text{mais} \qquad q \notin W, \]
\hfill q.e.d.

\msp{page 17}\textbf{Recollement dans} Cat \textbf{d'un ouvert et d'un fermé complémentaire.}

\texttt{[Supposons donné]}
\[ X_0 \underbrace{\hookrightarrow}_{\text{ouvert}} X \ \underbrace{\hookleftarrow}_{\substack{\text{fermé}\\ \text{complé-}\\ \text{mentaire}}} X_1 \]
dans Cat, \texttt{[et soit]}
\[ \tilde{X}_{01} \subseteq \mathrm{Fl}(X) \]
\texttt{[la]} sous-catégorie pleine des flèches de source dans $X_0$, but dans $X_1$.

\textbf{Question.} \textit{Ce diagramme}
\[
\begin{tikzcd}[column sep=small]
 & & \tilde{X}_{01} \arrow[dll, "p"' near start, "\text{fibrant, fibres les $x_0\backslash X_1$}"' near end] \arrow[drr, "q" near start, "\text{cofibrant, fibres les $X_0/x_1$}" near end] & & \\
X_0 \arrow[ddrr, "i"'] & & & & X_1 \arrow[ddll, "j"] \\
 & {} \arrow[rr, bend left] & & {} & \\
 & & X, & &
\end{tikzcd}
\tag*{$(*)$}
\]
\printedpage{13}
% head: A. Grothendieck : Dérivateurs, Ch. XII, Édition des Universités de Montpellier II et Paris VII
% foot: Version 20 juin 2003
\textit{est-il $W$-cocartésien?}

OK si $i$ $W$-asphérique ou j $W$-coasphérique, i.e. $X \longrightarrow \boldsymbol{\Delta}^1$ lisse ou propre -- c'est un cas où on a soit $X_0 \xrightarrow{W} X$, soit $X_1 \xrightarrow{W} X$. Mais même en supposant $i$ ou $j$ dans $W$, et de plus $X$ asphérique (donc $X_0$ et $X$ asphériques, ou $X_1$ et $X$ asphériques), et $W = \mathbf{W}_\infty$, je n'ai pu prouver que le diagramme est cocartésien, i.e. $q \in W$ resp. $p \in W$ (\footnote[9]{Je suppose que $W$ satisfait Loc (6).}). (Pas non plus si $X_0 \simeq e$, $X$, $X_1$ $W$-asphérique $\Longrightarrow$ $e\backslash X$ $W$-asphérique, i.e. $X_1 \longrightarrow X$ $W$-coasphérique $\ldots$\,) Mais j'ai prouvé du moins, dans tous les cas, que % [sic] j (upright) -> $j$, letter denoting a functor printed in roman type
\[
\begin{tikzcd}[row sep=small]
 & X_0 & & \\
\mathcal{X} \overset{\text{déf}}{=} \displaystyle\int \tilde{X}_{01} \arrow[ur] \arrow[dr] & & {} \arrow[r, "\varphi"] & X \\
 & X_1 & &
\end{tikzcd}
\tag*{$(**)$} % ?? uncertain: script letter X -- \mathcal{X} or \mathscr{X}
\]
est dans $\mathbf{W}_1$.

Chose \textit{fausse} (sauf si $i$ ou $j$ dans $W = \mathbf{W}_\infty$) si on remplace $\tilde{X}_{01}$ par $\tilde{X}'_{01}$, formée des triples $(x_0, x_1, c)$, avec $x_0 \in \mathrm{Ob}\,X_0$, $x_1 \in \mathrm{Ob}\,X_1$, $c : x_0 \dotto x_1$ un \textit{chemin} de $x_0$ à $x_1$ dans $X$! On a
\msp{page 18}donc une situation ‘topologique’ tout à fait spéciale à la catégorie de modèles Cat, qui ne semble pas avoir d'analogue dans des catégories telles que Esp ou $\boldsymbol{\Delta}^\wedge$, où pourtant les notions d'ouvert et de fermé complémentaire (avec grain de sel dans $\boldsymbol{\Delta}^\wedge$) ont un sens.

L'argument catégorique quasi-trivial qui établit que $(**)$ est dans $\mathbf{W}_1$, i.e. induit une équivalence de catégories entre $\mathrm{Rev}(X)$ et
\[ \mathrm{Rev}(\mathcal{X}) \ \simeq \ \mathrm{Rev}(X_0) \overset{2}{\times}_{\mathrm{Rev}(\tilde{X}_{01})} \mathrm{Rev}(X_1), \]
me semble devoir s'étendre aux dimensions supérieures, à coups de $n$-gerbes sur $X$ -- mais les fondements sur la théorie des $n$-gerbes manquent. (Et je n'ai pas eu le courage d'expliciter l'argument même pour $n = 2$ (où pourtant \textsc{Giraud} fournit ce qu'il faut) (\footnote[10]{Mais j'ai fait quand même une esquisse de vérification pour $n = 2$, qui me convainc qu'on a bien $\varphi \in \mathbf{W}_2$. Et je n'ai pas guère de doute qu'on a même $\varphi \in \mathbf{W}_\infty$.}).) Aussi je conjecture que $\varphi$ est tout au moins dans $\mathbf{W}_\infty$, i.e. que le diagramme $(*)$ est $W$-cocartésien si $W \supseteq \mathbf{W}_\infty$ (et satisfait Loc (6 a$'$)). % [sic] je n'ai pas guère (footnote 10) -> je n'ai guère, double negation pas + guère

\end{document}
